\chapter*{Preface}
\label{chap:preface}

This book is about a compiler that does not know it is doing physics.

The compiler in question is Lean 4, a proof assistant developed at Microsoft
Research. Lean's job is straightforward: given a definition, check whether it
is consistent with everything that has been defined before. Given an inference,
verify that it follows. Given a program, decide whether it compiles. The
compiler makes these decisions one at a time, in order, without looking ahead.
It cannot accept a new term until it has closed the terms that came before.

This book argues that this process -- the incremental, ordered, irrevocable
accumulation of verified decisions -- is the same process that produces
physical law.

That is not a metaphor. The claim is precise. A measurement is a decision made
by an instrument: this event is distinguishable from that one. A ledger is the
ordered, irreversible record of those decisions. A physical law is a pattern
that persists across every consistent extension of the ledger. The compiler's
elaborator closes a term in exactly the way a law closes a ledger: by
confirming that the new entry is consistent with everything recorded so far.
The compiler is a measuring instrument. The compilation is the measurement.

The book is organized as a guided tour through a Lean formalization called
\textit{the device}. The device was built in fifteen episodes, each episode
adding a new layer to a growing tower of type classes. By the final episode,
the tower has thirty-two layers, and the compiler is asked to verify a single
proposition: that the entire structure closes on its own truthhood. The
\lean{\#check} command that terminates the device is not a curiosity. It is the
experiment. The compiler's answer is the result.

A companion volume, \textit{Measurement: The Finite Space-Time Geometry of the
Single Invariant}, develops the same argument in full mathematical rigor. This
book is not that book. Here the goal is to make the ideas accessible to any
reader who has asked a compiler a question and wondered how it knew the answer.
The Lean code appears when it is the clearest way to say something. The physics
appears when it is the inevitable consequence of what came before.

Douglas Adams said the answer is 42. This book asks: what is the question, and
why does answering it require the entire structure of the universe?


